\section{Introduction: why a model needs an audit}
A model is a set of equations that has been put in a computer. Whether its results can be trusted is not a matter of how many equations it has but of what has been checked. An \kw{audit} is a set of such checks. The simplest and the hardest is the closure of the mass balance: what enters the root zone must equal what leaves it plus the change of storage. A model whose balance does not close is losing or creating water, and any saving computed from it is partly an artefact. After closure, the model is compared with data, the comparison is summarised by a few statistics, the uncertainty of the summary is stated, and the limits of what was done are written down. This chapter gives the tools for each of these steps.

\section{Verification, calibration and validation}
Three words are often mixed up. \kw{Verification} asks whether the equations have been solved correctly (is the code right? does it conserve mass? does it reproduce an analytic solution?). \kw{Calibration} adjusts uncertain parameters so that the model fits a set of observations. \kw{Validation} compares the calibrated model with observations \emph{not} used in the calibration. A model that has been verified and calibrated but not validated is a fitted curve; one that has been validated against data from a different site and season is a tested model. In this book the Daisy runs are verified by the balance audit (Chapter~\ref{chap:04}), the application's engines are verified against Daisy (Chapter~9), and no field validation is claimed.

\section{Goodness-of-fit statistics}
Let $O_i$ be $n$ observations (or the reference results) with mean $\bar O$ and $P_i$ the corresponding model values. The commonly used measures are the \kw{bias} $\bar e=\frac1n\sum(P_i-O_i)$, the \kw{mean absolute error} $\mathrm{MAE}=\frac1n\sum|P_i-O_i|$, the \kw{root-mean-square error} $\mathrm{RMSE}=\sqrt{\frac1n\sum(P_i-O_i)^2}$, the \kw{coefficient of determination} $R^2$ of the regression of $P$ on $O$, the \kw{Nash--Sutcliffe efficiency}

\begin{equation}\mathrm{NSE}=1-\frac{\sum_i(O_i-P_i)^2}{\sum_i(O_i-\bar O)^2},\label{eq:nse}\end{equation}
\par\noindent where NSE is the Nash--Sutcliffe efficiency (dimensionless, 1 for a perfect model), $O_i$ the observed and $P_i$ the predicted value at time $i$ and $\bar O$ the mean of the observations.



which is 1 for a perfect model and 0 when the model is no better than the mean of the observations, and the \kw{index of agreement} $d=1-\sum(P_i-O_i)^2/\sum(|P_i-\bar O|+|O_i-\bar O|)^2$ of Willmott, bounded between 0 and 1. No single measure is enough: $R^2$ ignores bias, RMSE is in the units of the data but hides the pattern, and NSE punishes outliers.

\exbox{Statistics of a model against data}{A model predicts the annual evapotranspiration (mm) of eight years: observed 312, 355, 298, 401, 377, 330, 289, 362; predicted 320, 349, 305, 390, 385, 338, 280, 371. Compute the bias, MAE, RMSE, $R^2$, NSE and $d$.}\sol Errors $P-O$: 8, -6, 7, -11, 8, 8, -9, 9. Bias $=1.75$ mm; MAE $=8.25$ mm; RMSE $=\sqrt{70.00}=8.37$ mm. Mean of $O$ is $340.50$; $\sum(O-\bar O)^2=11046$ and $\sum(O-P)^2=560$, so \[\mathrm{NSE}=1-560/11046=0.949\]
Correlation gives $R^2=0.952$; $d=0.987$ \ans The model is nearly unbiased (under 1 mm) and its error is about 2.5 per cent of the mean; NSE of 0.97 says that it explains almost all the variance of the eight years. \finans

\section{Closing the water balance}
For a run of $y$ years the balance of Chapter~2 gives, for each year, the residual $R=P+I-ET-D-Q-\Delta S$ (with $ET=T+E_s+E_o$). A model that conserves mass has $R=0$ to numerical precision. Where $R\ne0$ the audit classifies the cause: a flux that is not reported, a storage that is not counted (snow, ponds, litter, the mulch layer) or a numerical error. The \emph{unattributed evaporative term} of the studies is such a residual that is systematically positive in the mulch experiments: it is evaporation from the mulch layer that the standard output does not count. The balance is the basis of the claim that a saving in one flux is a gain in another, so it is audited first.

\section{Uncertainty of a mean and of a slope}
A mean of $n$ values with standard deviation $s$ has a standard error $s/\sqrt n$ and a 95 per cent confidence interval $\bar x\pm t_{0.975,n-1}s/\sqrt n$. A regression slope $b$ of $y$ on $x$ has standard error $s_b=s_e/\sqrt{\sum(x-\bar x)^2}$ with $s_e$ the residual standard deviation, and the interval $b\pm t\,s_b$. When the distribution is not normal or the statistic is awkward, the \kw{bootstrap} resamples the data with replacement many times and reads the interval from the percentiles of the statistic. An ensemble of simulations has no sampling noise in the usual sense (the same inputs give the same output), so the intervals describe the spread over soils and years, not experimental error; this is stated in every table that quotes them.

\exbox{Slope and its interval}{The data $(x,y)$ are $(1,2.1),(2,3.9),(3,6.2),(4,7.8),(5,10.1),(6,12.2),(7,13.8),(8,16.1)$. Fit a line, find the standard error of the slope and its 95 \% interval ($t_{0.975,6}=2.447$).}\sol \givenline{$t_{0.975,6}=2.447$}\ans $\bar x=4.5$, $\sum(x-\bar x)^2=42$. Slope $b=1.9976$, intercept $a=0.0357$. Residual standard deviation \[s_e=\sqrt{\sum e^2/6}=0.1802\]
\[s_b=0.1802/\sqrt{42}=0.0278\]
The interval is \[1.9976\pm2.447\times0.0278=[1.930,\,2.066]\quad\text{\finans}\]

\section{What a study of this kind can and cannot claim}
The findings of the studies concern a one-dimensional, matrix-flow model with ten soils, a fixed base case of crop and weather, and a designed set of changes. They say what the model does, with the cause separated from the others by the design; they do not say what a field does. Three limits follow. Matrix flow excludes cracks and macropores. One dimension excludes lateral flow and the heterogeneity of a field. A fixed base case limits the extrapolation to other crops and climates. A claim about the \emph{mechanism} (a saving of evaporation is not a saving of water, because the water is conserved) holds for any soil water model that conserves mass; a claim about the \emph{size} of an effect holds only for the model.

